% ============================================================
% Section 1: Introduction
% Structure follows CARS 5-Move model (NCA_WRITING_DNA.md)
% Citation policy: Prioritizes Neural Computing and Applications (NCA) literature
% ============================================================

\section{Introduction}\label{sec:introduction}

% --- Move 1: Territory (Domain motivation, market specifics, challenges) ---


Financial time-series forecasting remains a challenging computational problem because market data exhibit non-stationarity, nonlinear dependence, structural changes, and substantial noise \cite{jesus2025tda, makridakis2018statistical}. These characteristics become particularly relevant in short-horizon trading, where predictive relationships may weaken or change as market conditions evolve. Among traded financial instruments, futures contracts provide an important setting for this problem because they are widely used for both risk hedging and speculative exposure \cite{jarrow1981forward, souza2025iccsa}. In Brazil, the Mini-Index futures contract (symbol: WIN), traded on Brasil Bolsa Balc\~{a}o (B3) and referenced to the Ibovespa index, is highly liquid and extensively used in intraday operations, making it a relevant environment for computational financial research \cite{souza2024dataset, cardoso2022bovdb}. Predicting directional movements in such markets is therefore scientifically demanding and requires methods capable of extracting useful information from noisy and temporally evolving observations.

% --- Move 2: Model & Search Space (ML models, features, feature selection) ---
Traditional financial forecasting has relied on statistical models such as autoregressive integrated moving average (ARIMA) techniques \cite{box2015time}. As finance and computer science increasingly converged under the computational intelligence paradigm \cite{jiang2021applications}, machine-learning models became widely investigated for financial prediction and classification. Supervised approaches---including Random Forests (RF) \cite{nabipour2020predicting}, Support Vector Machines (SVM) \cite{naik2019optimal}, Multilayer Perceptrons (MLP) \cite{ecer2020mlp_ga_pso}, and 1D Convolutional Neural Networks (1D-CNN) \cite{bandgar2025finchronos}---can capture nonlinear relationships between market variables and future price movements. These models are commonly supplied with technical indicators derived from price, volume, momentum, and volatility. However, large indicator sets may introduce redundant or weakly informative variables, increasing model complexity and potentially degrading generalization \cite{singh2022feature, lei2012feature}. Feature-selection techniques therefore provide an important mechanism for identifying compact and informative predictor subsets \cite{kumar2014feature, souza2025iccsa}.

% --- Move 3: Optimization Paradigm (Hyperparameter tuning, MOAs) ---
Beyond feature selection, predictive performance also depends strongly on hyperparameter configuration \cite{jovanovic2024wann}. Neural and machine-learning models contain structural, regularization, and training parameters that jointly define mixed discrete--continuous and often non-convex search spaces. Gradient-based learning algorithms optimize model parameters once an architecture is defined, but they do not directly solve this outer hyperparameter-search problem when discrete and non-differentiable choices are involved \cite{rajwar2023exhaustive_review}. Metaheuristic Optimization Algorithms (MOAs) provide an alternative by exploring candidate configurations through stochastic search. Evolutionary approaches such as Genetic Algorithms (GA) \cite{goldberg1989ga} and Differential Evolution (DE) \cite{storn1997de}, together with swarm-based methods such as Particle Swarm Optimization (PSO) \cite{kennedy1995pso} and Grey Wolf Optimizer (GWO) \cite{faris2018gwo_review}, have consequently been applied to automated model tuning across neural computing problems \cite{al_bannoud2026accelerating, bashir2026evobagnet}.

% --- Move 4: Literature Gap (3 methodological limitations) ---
Despite the increasing use of metaheuristic optimization in predictive modelling, comparisons between algorithms remain strongly influenced by experimental design. Three methodological limitations are particularly relevant in financial time-series applications:

\begin{enumerate}
    \item \textbf{Unequal computational budgets}: Metaheuristics are commonly expressed using different units such as generations, iterations, populations, particles, or wolves. As a result, apparently similar experimental configurations may correspond to substantially different numbers of objective-function evaluations. Without a common evaluation budget, observed differences may reflect unequal search opportunities rather than differences in optimizer behaviour.

    \item \textbf{Temporal data leakage}: Randomized data partitioning and conventional shuffled cross-validation can violate the chronological structure of financial observations \cite{makridakis2018statistical}. When future observations contribute directly or indirectly to model selection, validation results may provide an overly optimistic estimate of out-of-sample generalization.

    \item \textbf{Dependence on the optimization objective}: Hyperparameter-search outcomes depend not only on the optimizer but also on the metric used to rank candidate configurations. Accuracy alone may provide an incomplete assessment of classification quality, particularly when prediction errors or class distributions are unbalanced \cite{chicco2020mcc}. Moreover, improvements in predictive metrics do not necessarily translate into economically useful trading behaviour once transaction costs and signal frequency are considered.
\end{enumerate}

% --- Move 5: Solution, Framework & Contributions ---
To address these limitations, this paper presents a controlled multi-optimizer benchmark for intraday trend classification using 5-minute B3 Mini-Index futures data. A compact seven-indicator subset previously identified through Information Gain is used as a fixed input representation, while five optimization strategies---Random Search (RS), Genetic Algorithm (GA), Particle Swarm Optimization (PSO), Differential Evolution (DE), and Grey Wolf Optimizer (GWO)---are evaluated across four classifier families: MLP, SVM, RF, and 1D-CNN. All optimizers operate under the same objective-function evaluation budget of $N_{\mathrm{eval}} = 1{,}000$ evaluations per seed, and temporal ordering is preserved through chronological validation with an isolated final test block. The study further contrasts a compound MCC-driven objective ($0.60\,\mathrm{MCC} + 0.40\,F_1$) with accuracy-based optimization protocols. Predictive performance is complemented by convergence analysis, non-parametric statistical testing, and out-of-sample economic evaluation under transaction costs.

The main contribution of this study is the development of a controlled experimental framework that disentangles the effects of optimizer family, predictive model, validation objective, and computational budget in intraday financial classification. Rather than evaluating metaheuristics under isolated or algorithm-specific settings, the proposed benchmark places five optimization strategies and four predictive model families under the same temporal data partitions, the same feature representation, and an equal number of objective-function evaluations. By combining predictive metrics, convergence behaviour, matched statistical testing, and economic evaluation under transaction costs, the study provides a more complete assessment of optimizer robustness and practical relevance. In this sense, the proposed framework extends existing approaches to intraday futures classification by integrating feature selection, multiple optimization strategies, temporal validation, and equal computational budgets within a unified and reproducible comparative methodology.

% --- Move 6: Outline ---
The remainder of this manuscript is organized as follows. Section~\ref{sec:literature_review} presents the related literature and establishes the research gap. Section~\ref{sec:methods} describes the dataset, feature selection, predictive models, optimization algorithms, search spaces, and experimental protocol. Section~\ref{sec:results} reports the predictive results, convergence behaviour, statistical analyses, and economic evaluation. Section~\ref{sec:discussion} discusses the findings, methodological implications, and limitations. Finally, Section~\ref{sec:conclusion} summarizes the main conclusions and outlines future research directions.